\section{The integral circle quotient and its cusp replacement}\hypertarget{tag-02OO}{}\label{mapsI:circle-quotient}

\begin{definition}[Circle-quotient conventions]\hypertarget{tag-IP06}{}\label{I:def:circle-quotient-conventions}
We write $g,u,w,d$ for the ordered homology basis
$(\widehat\gamma,\widehat u,\widehat w,\widehat\delta)$ in this
section, and $\gamma$ for its first dual covector. The cusp coordinate is
$q=1/t$.  Positive target orientations are the complex orientation
on the base followed by the positive $\gamma$-circle orientation.
The complex orientation on $X$ fixes the source orientation.
\end{definition}

\begin{proposition}[The primitive integral quotient and its finite fillings]\hypertarget{tag-02PA}{}\label{mapsI:src:A.1.1}
\StatementRecord{mapsI:src:A.1.1}{proved}{02PA}{\cite{AC2026}}

Use the monodromy matrices of \eqref{wb:pshift:eq:finite-data} and the
standard affine vectors of Definition~\ref{I:def:family-parameter}. Put
\[
 M_0=(A_1A_2)^{-1}=
 \begin{pmatrix}1&0&0&0\\0&1&0&0\\0&1&1&0\\-1&0&0&1\end{pmatrix}.
\]
The common invariant vectors and covectors are $\mathbb Zd$ and
$\mathbb Z\gamma$, respectively.  The smallest invariant primitive
sublattice containing the cusp vanishing lattice is
$L=\mathbb Z\langle u,w,d\rangle=\ker\gamma$.
The formula
\[
 [z,\Pi(z)x]\longmapsto([z],\gamma(x)\bmod\mathbb Z)
 \label{mapsI:quotient-eq-1}
\]
defines a real-analytic fibrewise group epimorphism
$J\longrightarrow B^\circ\times S^1$ with connected kernel $T_L$.
At the finite fillings it descends to a smooth proper $T^3$-bundle
over a Seifert solid torus, of multiplicities $3$ and $4$ and
translation residues $1$ and $-1$.
\end{proposition}



\begin{proof}
\begin{enumerate}[label=\arabic*.,leftmargin=*]
\item For the matrices of \eqref{wb:pshift:eq:finite-data}, direct integral
multiplication gives
$A_1^3=A_2^4=A_1A_2M_0=I$ and $A_jv_j=v_j$.
Solving the two fixed-vector equations gives $(0,0,0,d_0)$.
For a row vector $(a,b,c,e)$, the first fixed-covector equation
gives $b=2c,e=3c$ and the second gives $b=c,e=2c$;
therefore $b=c=e=0$.  Also
\[
 (M_0-I)(a,b,c,e)=(0,0,b,-a),\qquad
 A_1w=u-w.
\]
These identities prove the assertions about $L$ and the cusp
lattice $K=\mathbb Z\langle w,d\rangle$.  Their primitivity follows
from the displayed bases.  The flat real marking identifies the
deck action on a fibre with $x\mapsto A_jx$.  Since
$\gamma A_j=\gamma$, \eqref{mapsI:quotient-eq-1} descends.  Its kernel is connected
because $\gamma$ is primitive.

\item On a root disc the affine generator is
\[
 (s,x)\longmapsto
 (e^{-2\pi i/m}s,A_jx+v_j/m),\qquad (m,\kappa)=(3,1),(4,-1).
\]
It induces $(s,\theta)\mapsto
(e^{-2\pi i/m}s,\theta+\kappa/m)$ on the character quotient.
This action is free since $\gcd(\kappa,m)=1$.  Writing
$a=e^{2\pi i\theta}$ and choosing $r$ with $r\kappa\equiv1\pmod m$,
the invariant coordinates $s a^r$ and $a^m$ identify its quotient
with $D^2\times S^1$.  The base coordinate becomes
$(s a^r)^m(a^m)^{-r}$, and hence has multiplicity $m$ at its
central circle.  Before the finite quotient the map is a trivial
$T_L$-bundle.  The free actions on source and target are
equivariant for that map, so it descends locally to a smooth
$T_L$-bundle.  Properness follows from compactness of $T_L$.
The logarithmic translation used to glue the root chart to $J$
changes the character by its $\gamma$-component.  Applying
$\gamma$ to the actual gluing translation therefore gives the
target gluing, proving agreement on every overlap.
\end{enumerate}
\end{proof}


\begin{remark}[Comparison of the quotient characters]\hypertarget{tag-IP08}{}\label{I:rem:circle-character-comparison}
Engel's primitive invariant covector $\psi$
\cite[Lemma~3.8]{Engel2026} is related to the present character by
\[
 P_E=P_{\rm mark},\qquad \psi P_E=-\gamma,
 \hypertarget{tag-02OZ}{}\label{mapsI:quotient-eq-2}
\]
where $P_{\rm mark}$ is the integral isomorphism of
Proposition~\ref{inv:B1.0}. Thus the quotient characters agree after the
target-circle map $\theta\mapsto-\theta$. Inverting the root generators
gives the translation vectors $a_{j-1}=-P_Ev_j/m_j$, as in
Proposition~\ref{fam:src:B.2.2}.
\end{remark}

\begin{proposition}[Pointwise cusp limits and the fixed double curve]\hypertarget{tag-02PB}{}\label{mapsI:src:A.1.2}
\StatementRecord{mapsI:src:A.1.2}{proved}{02PB}{\SeriesRef{mapsI:quotient-eq-4}}

The literal character has no continuous extension at either
triple point or at generic points of two of the three double
curves.  It has the constant limit $0$ on the smooth locus of
$W$ and on the third double curve away from the triple points.
In the chart $q=z_0z_1z_2$, $x_1=z_1$, $x_2=z_2$, the extendible
double curve is $z_0=z_1=0$, $z_2\ne0$.  The fixed curve of the
$\mathbb C^*$-action in direction $d=e_2$ is instead
$z_0=z_2=0$, $z_1\ne0$, and is one of the nonextendible curves.
Consequently the extendible double curve and the fixed double curve are distinct.
\end{proposition}



\begin{proof}
\begin{enumerate}[label=\arabic*.,leftmargin=*]
\item Put $q=e^{-2\pi T}$ on a positive radial ray and use the exact
cusp periods
\[
 Z_T=iTB_0+C(q),\qquad
 B_0=\begin{pmatrix}0&1\\-1&0\end{pmatrix},
\]
where $C$ is bounded and holomorphic at $0$.  For a fixed
$a\in\mathbb R^2$ the representative
$x_T(a)=\exp(2\pi iZ_Ta)$ has character $a_1\bmod\mathbb Z$ and
logarithmic position
\[
 \frac{\log|x_T(a)|}{\log|q|}
 =B_0a+T^{-1}\operatorname{Im}C(q)a.
 \hypertarget{tag-02P0}{}\label{mapsI:quotient-eq-3}
\]
The choices $a=(-1/4,1/4)$ and $a'=(-1/2,1/4)$ have logarithmic
positions tending to $(1/4,1/4)$ and $(1/4,1/2)$.
Both are interior to the triangle with vertices $0,e_1,e_2$,
so all three chart coordinates tend to $0$.  Their character
values are $3/4$ and $1/2$.  This proves nonextension at the
first triple point.  The other triangle type has vertices
$e_1,e_2,e_1+e_2$.  Its coordinates can be taken as
$x_1x_2/q,q/x_2,q/x_1$.  Positions $(3/4,3/4)$ and
$(3/4,1/2)$ lie in its interior, give distinct values of $-y_2$
modulo one, and converge to its origin.  Translation by the
lattice preserves triangle type; these are precisely the two
triple points of $W$.

\item For arbitrary convergent paths in the first chart, inversion
of the exact logarithmic period matrix gives
\[
 \theta_\gamma=-\frac{\log|z_2|}{\log|q|}+o(1)\pmod{\mathbb Z}.
 \hypertarget{tag-02P1}{}\label{mapsI:quotient-eq-4}
\]
Here the error tends to zero whenever the logarithmic ratios
stay bounded; they do stay bounded near a fixed point of a
normal-crossing stratum because $q=z_0z_1z_2$ and the three
coordinates are bounded.  If $z_0,z_1\to0$ and $z_2\to c\ne0$,
the right side tends to zero.  If $z_0,z_2\to0$ and $z_1=c\ne0$,
the paths $z_2=q^a,z_0=q^{1-a}/c$, $0<a<1$, have limits
$-a\pmod1$ at the same double point.  The same construction
with $z_0=c$ proves nonextension on $z_1=z_2=0$.
At a smooth point of a component the ratio in \eqref{mapsI:quotient-eq-4} tends
to $0$ or $1$, giving zero modulo one.  Integer translates
of these charts give the stated assertion on all of $W$.

\item The torus action in direction $e_2=d$ is
\[
 \lambda\cdot(z_0,z_1,z_2)
       =(\lambda^{-1}z_0,z_1,\lambda z_2).
\]
Its fixed curve is $z_0=z_2=0$, while the character extends along
$z_0=z_1=0$ by the preceding limit formula. The displayed action and
limit formula give these two loci in the same toric chart.
\end{enumerate}
\end{proof}

\begin{lemma}[The cusp neighbourhood and its circle classes]\hypertarget{tag-IP07}{}\label{I:lem:cusp-fundamental-group}
For the cusp neighbourhood $N_0$ and the vanishing lattice
$K=\mathbb Z\langle w,d\rangle$,
\[
 \pi_1(N_0)=\overline\Lambda=\Lambda/K.
 \hypertarget{tag-02P2}{}\label{mapsI:quotient-eq-5}
\]
The punctured-cusp map to this group kills $K$ and the base meridian.
Both long-period covectors extend in integral $H^1(N_0;\mathbb Z)$.
The literal character on an ordinary fibre has no extension to all of
$X\simeq S^6$.
\end{lemma}
\begin{proof}
The toric covering has rays $(v,1)$ for all $v\in\mathbb Z^2$.
Meridians about its toric divisors kill the whole cocharacter lattice:
$(0,1)$ is a ray and $(e_i,1)-(0,1)=(e_i,0)$. Hence the toric cover is
simply connected. Its free period deck group is
$\overline\Lambda=\Lambda/K$, proving the displayed isomorphism.
The base meridian bounds the disc given by a point of the dense orbit
on a central component. Both long-period covectors kill these
relations and therefore extend in integral $H^1$.
An extension of $\theta_\gamma$ to all of $X\simeq S^6$ would give a
nonzero homomorphism $H_1(F)\to\mathbb Z$ factoring through
$H_1(X)=0$, which is impossible.
\end{proof}

\begin{proposition}[The filled three-dimensional target]\hypertarget{tag-02PC}{}\label{mapsI:src:A.1.3}
\StatementRecord{mapsI:src:A.1.3}{proved}{02PC}{the definitions and constructions stated in the proof}

For the original data the quotient over the complement of
the cusp disc is the exterior of the $(3,4)$ torus knot,
with the specified Seifert structure.  For general admissible
$(l_0,l_1,l_2)$ the filled quotient is the lens space determined
by the following meridian slopes in one oriented integral
torus basis $(\rho_1,h)$:
\[
 \mu_1=(3,-l_1),\qquad
 \mu_2=(-4,4l_0-l_2),\qquad
 \det(\mu_1,\mu_2)=p=12l_0-4l_1-3l_2.
 \label{mapsI:quotient-eq-6}
\]
In the lens-space convention that the second meridian is
$-q_0\mu_1+p\lambda_1$, where
$\det(\mu_1,\lambda_1)=1$, the lens parameter satisfies
$3q_0\equiv4\pmod p$.  In particular the target depends only
on the signed integer $p$, and is a three-sphere exactly
when $|p|=1$.
\end{proposition}



\begin{proof}
\begin{enumerate}[label=\arabic*.,leftmargin=*]
\item Let $P$ be the three-holed sphere, with oriented boundary
meridians $\rho_0+\rho_1+\rho_2=0$ and positive fibre $h$.
The two finite meridian relations are
$3\rho_1-l_1h=0$ and $4\rho_2-l_2h=0$; the cusp solid-torus
meridian is $\rho_0+l_0h$.  Eliminating $\rho_0$ gives
$\rho_2=l_0h-\rho_1$ on the annulus joining the finite solid
tori, which gives exactly \eqref{mapsI:quotient-eq-6}.  Since both slopes are
primitive, this is a genus-one Heegaard splitting.  Choose
$a,b\in\mathbb Z$ with $al_1+3b=1$ and put
$\lambda_1=(a,b)$.  Direct multiplication gives
\[
 \mu_2=-q_0\mu_1+p\lambda_1,
 \quad q_0=a(4l_0-l_2)+4b,
 \quad 3q_0-ap=4.
 \hypertarget{tag-02P3}{}\label{mapsI:quotient-eq-7}
\]
This proves the claimed lens parameter in a convention fixed
by the displayed oriented meridians.  The fundamental
group of this gluing is $\mathbb Z/|p|$.  Admissibility implies
$\gcd(p,6)=1$, so the modular inverse of $3$ used above exists.

\item For $(l_0,l_1,l_2)=(0,1,-1)$ the two exceptional-fibre data
are $(3,1),(4,-1)$.  They are realized by the circle action
\[
 s\cdot(z_1,z_2)=(e^{4is}z_1,e^{3is}z_2)
 \quad\text{on } S^3\subset\mathbb C^2.
\]
Its invariant map is $[z_1^3:z_2^4]$; the local circle
translation residues are $4^{-1}=1\pmod3$ and
$3^{-1}=-1\pmod4$.  Removing a regular orbit gives the
same two Seifert solid tori and their product annulus as
above. A regular orbit is the $(3,4)$ torus knot under the coordinate
exchange $(z_1,z_2)\mapsto(z_2,z_1)$ when the opposite ordering is used.
The displayed orbit map and meridian slopes identify the manifold
with its peripheral structure.
\end{enumerate}
\end{proof}

\subsection{The cusp replacement and its relative uniqueness}

\begin{theorem}[The cusp replacement and its relative homotopy class]\hypertarget{tag-02PD}{}\label{mapsI:src:A.2.1}
\StatementRecord{mapsI:src:A.2.1}{proved}{02PD}{\SeriesRef{mapsI:quotient-eq-5}}

For every admissible member, there is a smooth map
$e_{\rm q}:X\to Q$ over $B$ whose restriction outside an arbitrarily
small cusp neighbourhood is the literal circle quotient.
For fixed quotient gluing, its homotopy class is independent
of the cusp replacement, provided the replacement maps the
cusp neighbourhood into its designated solid torus and is
kept fixed on its outer collar.  In particular there is a
well-defined class for the original member $X\simeq S^6$.
\end{theorem}



\begin{proof}
\begin{enumerate}[label=\arabic*.,leftmargin=*]
\item On the cusp boundary use the product coordinates
$(q,\theta_\gamma)$ after the specified collar shear.
By \eqref{mapsI:quotient-eq-5}, the boundary circle class extends to $N_0$.
Since $S^1=K(\mathbb Z,1)$, it has a continuous representative
$\theta':N_0\to S^1$.  Its restriction to the boundary is
homotopic to the given one: the two maps induce the same
homomorphism on $\pi_1$ of the boundary mapping torus.
The homotopy-extension property for a manifold with boundary
changes $\theta'$ on a collar so that it agrees exactly with
$\theta_\gamma$ there.  Relative smooth approximation gives
a smooth representative, still equal on a smaller collar.
Then $(f,\theta')$ is the required map into $D^2\times S^1$,
and agrees with the finite-filling maps on the overlap.
Both components of $(f,\theta')$ are smooth on $N_0$, so the
product map is smooth at the critical points of $f$ as well.

\item The inclusion $\partial N_0\to N_0$ surjects on $\pi_1$,
as \eqref{mapsI:quotient-eq-5} explicitly shows.  It follows that
$H^1(N_0;\mathbb Z)\to H^1(\partial N_0;\mathbb Z)$
is injective.  Since both spaces are connected, the long
exact sequence gives $H^1(N_0,\partial N_0;\mathbb Z)=0$.
Two circle replacements agreeing on the boundary therefore
have quotient whose relative $H^1$ class is zero, hence
have a relative homotopy.  Keeping the base coordinate
fixed gives the required homotopy of maps into the solid
torus.  Different choices of the collar size are compared
on a common smaller collar.  This proves the assertion.
\end{enumerate}
\end{proof}
